\section*{Práctica 2 --- ¿Quién podía hacer qué?}
\subsection*{Accesos, integridad y procesos con relevancia jurídica}

La organización continúa investigando el incidente analizado en la práctica anterior.

El equipo jurídico necesita ahora responder a varias preguntas:
\begin{itemize}
  \item ¿qué usuarios existen en el sistema?;
  \item ¿qué documentos podía consultar cada uno?;
  \item ¿estaban correctamente configurados los permisos?;
  \item ¿podemos comprobar si un archivo ha sido modificado?;
  \item ¿qué significa que un programa esté ejecutándose?
\end{itemize}

\textbf{Tiempo previsto: 35--40 minutos.} La práctica está diseñada para completarse en menos de 45 minutos.

\subsection*{1. ¿Quién es quién en el sistema?}

Ejecute:
\begin{verbatim}
whoami
id
id marta
id diego
id responsable
\end{verbatim}

Observe especialmente la pertenencia a los grupos \texttt{juridico} y \texttt{soporte}.

\paragraph{Pregunta.}
¿Qué usuarios pertenecen al grupo \texttt{juridico}? ¿Pertenece \texttt{diego} a dicho grupo?

\subsection*{2. ¿Quién puede leer cada documento?}

Examine los documentos:
\begin{verbatim}
ls -l /opt/caso02/documentos
\end{verbatim}

Los permisos tendrán un aspecto similar a:
\begin{verbatim}
-rw-r----- ... responsable juridico ... informe_interno.txt
-rw-r--r-- ... responsable juridico ... contrato_cliente.txt
\end{verbatim}

Compruebe los accesos:
\begin{verbatim}
su -c "cat /opt/caso02/documentos/informe_interno.txt" marta
su -c "cat /opt/caso02/documentos/informe_interno.txt" diego
su -c "cat /opt/caso02/documentos/contrato_cliente.txt" diego
\end{verbatim}

\paragraph{Pregunta.}
¿Por qué puede \texttt{diego} leer \texttt{contrato\_cliente.txt} si no pertenece al grupo \texttt{juridico}?

\subsection*{3. Corregir un permiso excesivo}

Modifique los permisos:
\begin{verbatim}
chmod 640 /opt/caso02/documentos/contrato_cliente.txt
ls -l /opt/caso02/documentos/contrato_cliente.txt
\end{verbatim}

Vuelva a comprobar:
\begin{verbatim}
su -c "cat /opt/caso02/documentos/contrato_cliente.txt" marta
su -c "cat /opt/caso02/documentos/contrato_cliente.txt" diego
\end{verbatim}

\paragraph{Pregunta.}
¿Permite este dato afirmar que \texttt{diego} nunca llegó a leer el contrato antes de corregir los permisos?

\paragraph{Respuesta.}
No. Los permisos describen una capacidad de acceso en un determinado momento. Por sí solos no demuestran si un acceso se produjo realmente.

\subsection*{4. Comprobar la integridad}

Calcule y guarde un hash SHA-256:
\begin{verbatim}
sha256sum /opt/caso02/documentos/contrato_cliente.txt
sha256sum /opt/caso02/documentos/contrato_cliente.txt \
> /opt/caso02/trabajo/hash_original.txt
cat /opt/caso02/trabajo/hash_original.txt
\end{verbatim}

Cree una copia y compare:
\begin{verbatim}
cp /opt/caso02/documentos/contrato_cliente.txt \
/opt/caso02/trabajo/contrato_cliente_copia.txt

sha256sum /opt/caso02/documentos/contrato_cliente.txt \
/opt/caso02/trabajo/contrato_cliente_copia.txt
\end{verbatim}

Modifique sólo la copia:
\begin{verbatim}
echo "Anotación añadida durante la revisión." \
>> /opt/caso02/trabajo/contrato_cliente_copia.txt
\end{verbatim}

Vuelva a comparar y verifique el original:
\begin{verbatim}
sha256sum /opt/caso02/documentos/contrato_cliente.txt \
/opt/caso02/trabajo/contrato_cliente_copia.txt
sha256sum -c /opt/caso02/trabajo/hash_original.txt
\end{verbatim}

Un hash coincidente proporciona una comprobación muy fuerte de integridad, pero no demuestra por sí solo quién creó o modificó un archivo.

\subsection*{5. De un programa a un proceso}

Sitúese en:
\begin{verbatim}
cd /opt/caso02/procesos
vi proceso.c
\end{verbatim}

e introduzca:
\begin{verbatim}
#include <stdio.h>
#include <unistd.h>

int main(void)
{
    printf("Proceso de prueba iniciado. PID=%d\n", (int)getpid());
    fflush(stdout);
    sleep(300);
    return 0;
}
\end{verbatim}

Compile y ejecute:
\begin{verbatim}
gcc proceso.c -o proceso
./proceso &
echo $! > proceso.pid
cat proceso.pid
\end{verbatim}

Un programa es un archivo ejecutable. Un proceso es una instancia de ese programa que se encuentra en ejecución.

\subsection*{6. Examinar y finalizar el proceso}

\begin{verbatim}
ps -p "$(cat proceso.pid)" -o pid,user,cmd
kill "$(cat proceso.pid)"
wait "$(cat proceso.pid)" 2>/dev/null
ps -p "$(cat proceso.pid)" -o pid,user,cmd
ls -l
\end{verbatim}

El proceso habrá desaparecido, pero el archivo ejecutable continúa existiendo.

\paragraph{Pregunta.}
¿Demuestra la salida actual de \texttt{ps} qué procesos estaban ejecutándose ayer?

\paragraph{Respuesta.}
No. \texttt{ps} muestra el estado de los procesos cuando se ejecuta la orden. Para estudiar actividad anterior serían necesarias otras fuentes de información.

\subsection*{7. Conclusiones}

Deben conservarse tres ideas:
\begin{enumerate}
  \item Tener permiso de acceso no demuestra que dicho acceso se haya producido realmente.
  \item Un hash permite detectar modificaciones del contenido, pero no demuestra autoría.
  \item Un proceso es una ejecución concreta de un programa y su observación actual no constituye por sí sola un historial.
\end{enumerate}

\appendix

\section*{Anexo A. Realización en Killercoda}

La práctica puede realizarse utilizando el escenario preparado por el profesor en Killercoda.

El entorno gratuito es temporal y la sesión tiene una duración máxima aproximada de una hora. Por ello, la práctica está diseñada para completarse en menos de 45 minutos. El sistema del escenario es efímero: al finalizar o reiniciar la sesión pueden perderse los cambios realizados.

Los botones \textbf{CHECK} comprueban determinados estados finales del sistema. No necesariamente permiten determinar qué secuencia exacta de órdenes utilizó el estudiante para alcanzar dicho estado.

\section*{Anexo B. Realización en un sistema Linux local}

Esta práctica crea usuarios y grupos de prueba. Se recomienda realizarla únicamente en una máquina destinada a prácticas, una máquina virtual o un sistema que pueda restaurarse posteriormente.

Cree un archivo denominado \texttt{preparar\_caso02.sh} con el mismo contenido que el archivo \texttt{background.sh} suministrado con la práctica y ejecútelo con:
\begin{verbatim}
sudo bash preparar_caso02.sh
sudo -i
\end{verbatim}

A partir de ese momento puede seguir los mismos pasos descritos en el cuerpo principal de la práctica.

El sistema necesita disponer al menos de:
\begin{verbatim}
vi
gcc
ps
sha256sum
su
\end{verbatim}
